\begin{table}[htbp]
\centering
\caption{Country-level decomposition of synthetic labor income}
\label{tab:latam_country_decomposition}
\small
\resizebox{\textwidth}{!}{%
\begin{tabular}{lrrrrrr}
\toprule
Economy & Initial & Final & Total change & Education & Within groups & Same source \\
\midrule
Argentina & 2016 & 2023 & -27.6 & 2.8 & -30.3 & Yes \\
Bolivia & 2016 & 2023 & -3.8 & 4.0 & -7.8 & Yes \\
Brazil & 2016 & 2023 & 5.6 & 9.3 & -3.7 & Yes \\
Chile & 2016 & 2023 & 10.1 & 17.5 & -7.4 & Yes \\
Colombia & 2016 & 2023 & 20.7 & 27.6 & -6.9 & No \\
Costa Rica & 2016 & 2024 & 0.5 & 7.3 & -6.9 & Yes \\
Dominican Republic & 2016 & 2023 & 4.5 & 2.7 & 1.9 & Yes \\
Ecuador & 2016 & 2024 & -17.2 & -3.1 & -14.2 & No \\
El Salvador & 2016 & 2023 & 22.9 & 3.1 & 19.8 & Yes \\
Guatemala & 2016 & 2022 & 3.3 & -2.5 & 5.8 & Yes \\
Mexico & 2016 & 2023 & 9.1 & 5.3 & 3.8 & Yes \\
Paraguay & 2017 & 2023 & -0.5 & 5.8 & -6.4 & No \\
Peru (Lima and Callao) & 2016 & 2023 & -10.8 & 2.9 & -13.6 & Yes \\
Uruguay & 2016 & 2023 & -0.1 & 4.2 & -4.3 & No \\
\bottomrule
\end{tabular}
}
\begin{minipage}{0.98\textwidth}
\footnotesize
\textit{Note:} Total change, education, and within-group components are percentages of initial synthetic monthly income. The underlying income measure is 2017 PPP US dollars per employed person per month. The two components add to total change before rounding. ``Same source'' indicates that the Tableau series and survey identifiers are unchanged across endpoints.
\end{minipage}
\end{table}
